\section{Research target and precise scope}\label{qspec:sec:context}

The canonical manuscript organizes polylogarithms, Gamma products, Hurwitz
spectral derivatives, and generalized Stieltjes constants within one exact
calculus \cite{repo}. Its current research chapter explicitly distinguishes
an analytic identity from a relation in a formal presentation and from a
small numerical residual. We retain that distinction. In particular, the
Gaussian $S_6$ identity is not proved here, and the rejection of one frozen
$S_8$ vector is not interpreted as rejection of every identity in its
comparison space.

The present continuation uses a different exact mechanism: introduce a
positive parameter $b$ into $x^2+b^2$, evaluate the resulting bilateral
spectrum in two ways, and then remove the deformation by a convergent
difference. The even spectrum sees the sum of reflected Hurwitz functions.
A first-moment spectrum sees their difference. Neither channel alone
recovers an individual Stieltjes constant, but together they do.

\subsection{The results and their logical dependence}
Theorem~\ref{qspec:thm:master} is the complete reflected Stieltjes identity.
Theorem~\ref{qspec:thm:polylog-integrals} replaces its Fourier--Bessel terms
by ordinary integrals of damped polylogarithms. The harmonic coefficient
formula and finite-head identity in Theorems~\ref{qspec:thm:harmonic}
and~\ref{qspec:thm:head} make all spectral derivatives explicit and avoid
an invalid extension of a binomial series outside its convergence disk.

The second main result, Theorem~\ref{qspec:thm:anomaly}, gives a complete
polynomial-weight parity criterion for the conjugate factorization defect.
It is independent of the Stieltjes master identity. Combining it with the
classical half-integer Bessel polynomials gives
Theorem~\ref{qspec:thm:finite}, a finite complex-Hurwitz/polylogarithm identity
at every negative integer. Theorem~\ref{qspec:thm:mass-ladder} identifies the
same family as repeated mass antiderivatives of a logarithm.

An illustrative member is obtained by setting
$z=\exp(-2\pi b+2\pi ia)$:
\begin{align}
&2\Rea\bigl\{\zeta'(-2,a+ib)+\zeta'(-2,1-a+ib)\notag\\[-2pt]
&\hspace{35mm}-2ib\bigl(\zeta'(-1,a+ib)+\zeta'(-1,1-a+ib)\bigr)\bigr\}\notag\\
&\qquad=-\frac{4\pi b^3}{3}
-\frac{2b}{\pi}\Rea\Li_2(z)-\frac1{\pi^2}\Rea\Li_3(z).
\label{qspec:eq:opening-example}
\end{align}
It holds for every $0<a<1$ and $b>0$. There is no unreported tail on the
right. Its proof is part of the all-index theorem, not an integer-relation
search.

\subsection{Attribution and what is not asserted}
The Gaussian Poisson identity, Mellin transforms, half-integer Bessel
functions, Hurwitz parameter differentiation, and Gamma reflection are
classical \cite{theta,bessel,halfK,hurwitz,gamma}. The first-derivative
multiplicative anomaly is also a classical subject; Dowker treats products
of linear factors and polynomial degeneracies \cite{dowker}. The repository
already has an arithmetic spectral-jet continuation that separates local
residue defects from nonlocal values \cite{arithmetic}. Accordingly, we do
not advertise the existence of the anomaly, Poisson continuation, or
negative-polygamma primitives as discoveries.

What is established here is an explicit two-channel Stieltjes construction,
its global harmonic implementation, a conjugate-pair polynomial parity
classification, and the resulting finite identities and primitive ladders.
These statements are derived in full below. Their absence from the portions
inspected is not an exhaustive literature or repository priority check.
No transcendence, numerical period independence, or proof-assistant
formalization is claimed.

\subsection{Conventions and analytic domains}\label{qspec:sec:conventions}
Unless stated otherwise,
\begin{equation}\label{qspec:eq:domain}
0<a<1,\qquad b>0,\qquad Q_b(x)=x^2+b^2,
\qquad x_n=n+a,\quad y_n=n+1-a.
\end{equation}
A one-sided formula needs only $a>0$. The variable $s$ is a spectral order;
$\zeta^{(r)}(s,a)$ and $\zeta'(s,a)$ differentiate that order, not $a$.
We use
\begin{equation}\label{qspec:eq:stieltjes-convention}
\zeta(1+u,a)=\frac1u+
\sum_{m=0}^{\infty}\frac{(-1)^m}{m!}\gamma_m(a)u^m,
\qquad G_m^\pm(a)=\gamma_m(a)\pm\gamma_m(1-a).
\end{equation}
In particular $\gamma_0(a)=-\psi(a)$. Superscripts on $\gamma_m^{(p)}$
will mean parameter derivatives.

All powers of positive $Q_b(x)$ use the real logarithm. For complex shifts
$a\pm ib$ we use the principal logarithm, and the branch of $\log\Gamma$
holomorphic on the right half-plane and real on its positive axis. Thus
$\log(x+ib)+\log(x-ib)=\log Q_b(x)$ for $x>0$, with no branch correction.
The damped polylogarithms have arguments strictly inside the unit disk,
where their defining series is entire in the order. Unit-circle limits
are stated separately.

The complete exponential Bell polynomials are defined by
\begin{equation}\label{qspec:eq:bell}
\exp\left(\sum_{r\ge1}v_r\frac{t^r}{r!}\right)
=\sum_{m\ge0}\Bcal_m(v_1,\ldots,v_m)\frac{t^m}{m!}.
\end{equation}
The convention $\Bcal_0=1$ and the recurrence
\begin{equation}\label{qspec:eq:bell-recurrence}
\Bcal_m=\sum_{r=1}^m\binom{m-1}{r-1}v_r\Bcal_{m-r}
\end{equation}
make every finite coefficient used below directly reproducible. We write
$H_k^{(r)}=\sum_{j=1}^k j^{-r}$ and $H_0^{(r)}=0$.
